\documentclass[11pt]{article}

\usepackage[utf8]{inputenc}
\usepackage[T1]{fontenc}
\usepackage{amsmath}
\usepackage[backend=bibtex,style=numeric,sorting=nyt]{biblatex}
\addbibresource{refs.bib}

\title{Evaluation Protocols for Retrieval-Augmented Generation: A Survey}
\author{A. Author \\ Department of Computer Science}
\date{May 2026}

\begin{document}

\maketitle

\begin{abstract}
Retrieval-augmented generation (RAG) systems are evaluated along at
least three largely independent axes --- retrieval quality, answer
faithfulness, and end-task accuracy --- and the literature has not
converged on which combination of metrics best predicts deployed
system quality. This survey organizes the evaluation protocols used
across twenty recent systems, groups them by what they actually
measure, and argues that faithfulness metrics in particular are
under-reported relative to how often unfaithful generations are cited
as a failure mode.
\end{abstract}

\section{Introduction}
\label{sec:intro}

Retrieval-augmented generation was popularized as a way to ground a
language model's output in retrieved evidence rather than parametric
memory alone \autocite{lewis2020rag}. Since then, evaluation practice
has fragmented: some papers report only end-task accuracy on a
question-answering benchmark \autocite{karpukhin2020dpr}, others report
retrieval metrics such as recall@$k$ in isolation
\autocite{izacard2021contriever}, and a smaller set reports
faithfulness --- whether the generated answer is actually supported by
the retrieved passages --- as its own axis \autocite{es2023ragas,niu2024ragtruth}.

\section{Retrieval Quality}
\label{sec:retrieval}

Retrieval quality is the most consistently reported axis, typically as
recall@$k$ or nDCG against a fixed passage index. \textcite{karpukhin2020dpr}
established dense passage retrieval as a strong baseline for open-domain
question answering, and \textcite{izacard2021contriever} showed that
contrastive pretraining on unlabeled text closes much of the gap to
supervised retrievers without requiring labeled query-passage pairs.

\subsection{What Recall@$k$ Does Not Capture}
\label{sec:recall-limits}

Recall@$k$ treats every relevant passage as equally useful once
retrieved, which is not true once the passage is handed to a generator
with a finite context window: a relevant passage buried at rank $k$
among $k{-}1$ irrelevant ones competes for attention with all of them.
\Autocite{liu2023lost} showed that generation quality degrades measurably
when the relevant passage sits in the middle of a long retrieved context,
even when recall@$k$ is unchanged --- a result that recall-only
evaluation protocols cannot detect at all.

\section{Faithfulness}
\label{sec:faithfulness}

Faithfulness --- whether the generated answer is actually entailed by
the retrieved evidence, as opposed to merely topically related to it ---
is reported far less consistently than retrieval recall. Two automated
protocols dominate recent work:

\begin{itemize}
  \item \textbf{Entailment-based scoring} \autocite[see][for the
    original formulation]{es2023ragas} uses a natural language inference
    model to check whether each claim in the generated answer is
    entailed by the retrieved passages.
  \item \textbf{Hallucination benchmarks} such as
    \autocite{niu2024ragtruth} construct held-out sets with known
    unsupported spans and report span-level precision and recall
    against human annotation.
\end{itemize}

Of the systems surveyed, fewer than half report a faithfulness metric
of either kind, despite unfaithful generation being cited as the
dominant failure mode in essentially every qualitative error analysis
we reviewed \autocite{niu2024ragtruth,es2023ragas}.

\section{End-Task Accuracy}
\label{sec:endtask}

End-task accuracy --- exact match or F1 against a reference answer on a
downstream benchmark --- is the most commonly reported number, and the
one most papers optimize against directly \autocite{karpukhin2020dpr,lewis2020rag}.
It is also the number most confounded by the other two axes: a system
can achieve high end-task accuracy with mediocre retrieval, if the
generator's parametric knowledge covers the gap, which is exactly the
behavior RAG is meant to reduce reliance on in the first place.

\section{Recommendations}
\label{sec:recs}

We recommend reporting all three axes rather than treating end-task
accuracy as a sufficient summary statistic, and in particular reporting
a faithfulness metric whenever a system claims to reduce hallucination
--- a claim that end-task accuracy alone cannot substantiate, and that
recall@$k$ cannot substantiate either, per the discussion in
Section~\ref{sec:recall-limits}.

\printbibliography

\end{document}
